% Methods section — draft template with our real numbers filled in.
% Written in simple words on purpose. Read every sentence and change it
% to your own voice where you can. Anything in [[double brackets]] is a
% spot where YOU must write the sentence yourself.

\section{Methods}

We simulate the manufacturing of a 40-layer mirror coating made of
silica (SiO$_2$) and titania-doped tantala (Ti:Ta$_2$O$_5$), the same
material pair used in Advanced LIGO's end mirrors. Every layer gets a
random thickness error, as happens in a real coating machine. We then
test a simple idea: after most of the layers are made, re-design the
few layers that are not made yet, so they cancel the errors in the
layers below. We call this step the \emph{rescue}. All of this is done
in simulation, on a laptop. No real coating is made.

The parts below describe the optical and noise models (Sec.~1.1), the
two mirror designs we test (Sec.~1.2), the error model (Sec.~1.3), the
rescue step (Sec.~1.4), and how we count a mirror as good or bad
(Sec.~1.5).

\subsection{Optical and thermal noise models}

We compute how much light a coating transmits with the transfer-matrix
method, at two wavelengths: 1064~nm (the main laser) and 532~nm (the
green locking beam), both at normal incidence. We use the open-source
\textsc{OptimalBragg} package from Caltech \cite{Venugopalan2024} for
these calculations. Material values are in Table~\ref{tab:materials}.

A warm coating also shakes a tiny bit on its own (Brownian thermal
noise), and this shaking is one of the main noise sources in LIGO. We
compute this noise following Hong \emph{et al.}~\cite{Hong2013}, as
implemented in \textsc{OptimalBragg}, with a laser beam radius of
6.2~cm (the Advanced LIGO end-mirror value~\cite{aLIGO2015}) and a temperature of 295~K. We report it as the amplitude
spectral density at 100~Hz. Inside optimization loops we use a faster
stand-in: the sum of each material's total thickness weighted by its
loss angle and stiffness, following LIGO document E0900068. The full
formula is used for all final numbers.

\begin{table}[h]
\centering
\caption{Material values used in all simulations, taken from the
\textsc{OptimalBragg} package \cite{Venugopalan2024}, which cites
arXiv:0912.0107~\cite{MaterialsSource}.}
\label{tab:materials}
\begin{tabular}{lcc}
\hline
 & SiO$_2$ & Ti:Ta$_2$O$_5$ \\
\hline
Refractive index (1064 nm) & 1.45 & 2.06 \\
Young's modulus (GPa) & 60 & 140 \\
Poisson ratio & 0.17 & 0.23 \\
Mechanical loss angle & $5\times10^{-5}$ & $2\times10^{-4}$ \\
\hline
\end{tabular}
\end{table}

Published values of the Ti:Ta$_2$O$_5$ loss angle vary; for example,
Ref.~\cite{Venugopalan2024} uses $3.6\times10^{-4}$. A larger value
raises every noise number by a similar factor (26--29\% for this
change), but it does not affect our results: pass rates depend only on
transmission, and every noise result we report is a comparison between
designs or between rescued and unrescued mirrors.

\subsection{The two mirror designs}

Both designs have 40 layers (20 pairs) plus a half-wave SiO$_2$ cap on
top, and both aim at the same two targets: a transmission of 5~ppm at
1064~nm and 2.5\% at 532~nm. We picked 2.5\% because it sits in the
middle of the allowed window (1--4\%, see Sec.~1.5), so errors have
room on both sides.

We made the designs with the differential evolution optimizer in
SciPy, on top of \textsc{OptimalBragg}'s optics and noise functions.
% [[One sentence in your own words on what the optimizer does, e.g.
% it keeps adjusting all 40 layer thicknesses until the coating hits
% both transmission targets with as little thermal noise as it can.]]

Running the optimizer from different random starting points gave
designs of two kinds, and we kept one of each:
\begin{itemize}
\item a \emph{quiet} design, with a Brownian noise of
$6.19\times10^{-21}\,\mathrm{m}/\sqrt{\mathrm{Hz}}$ at 100~Hz, which
turned out to be fragile against thickness errors;
\item a \emph{sturdy} design, about 8\% noisier
($6.68\times10^{-21}\,\mathrm{m}/\sqrt{\mathrm{Hz}}$), which tolerates
errors much better.
\end{itemize}
Testing both lets us ask whether the rescue can make the quiet design
usable, or only helps a design that is already sturdy.
As a check that our results do not hang on one lucky design, we also ran
the main tests on a third design from the same search, between the two in
robustness (5.12~ppm and 2.45\%).

\subsection{Deposition error model}

Layers are deposited one at a time, from the substrate up. The true
thickness of layer $i$ comes out as
\begin{equation}
d_i = d_i^{(0)}\,(1+\epsilon_i),
\end{equation}
where $d_i^{(0)}$ is the planned thickness and $\epsilon_i$ is a
random error drawn from a normal distribution with zero mean and
standard deviation $\sigma$. Each layer gets its own independent
error. We test $\sigma$ from 0.5\% to 3\%; around the interesting
region we use steps of 0.25\%.

For comparison, the Monte Carlo shipped with \textsc{OptimalBragg}
uses one shared error for all layers (the whole stack gets thicker or
thinner together). We also ran our test with that shared-error model.
It is much harsher: at a 1\% error the sturdy design passes 73\% of the
time (against 100\% with independent errors), and rescuing the top four
layers raises this to 92\%; at 2\% the rescue only lifts it from 42\% to
56\%. We report both models in the results.

\subsection{The rescue step}

In each simulated run, every layer except the top $k$ (and the cap) is
deposited with the random errors above. We then assume the machine has
measured the deposited layers, with its own error: each measured
thickness is the true one times $(1+\mu_i)$, where $\mu_i$ is normal
with standard deviation 0, 0.25, 0.5, 1 or 2\%. No published source
gives one agreed accuracy for in-situ monitoring, so we sweep this
value instead of picking one.

Using the measured stack, we re-optimize only the top $k$ layers
(the half-wave cap is never changed), aiming at the same two targets
with the same small noise penalty. We use the L-BFGS-B optimizer
started from three points (the planned thicknesses, and $\pm10\%$).
We checked it against the slower differential evolution on 1600
mirrors ($k=4$, both designs, both error models, 2--3\% errors, 0 and
1\% measurement error). For the sturdy design the two disagreed on 4 of
800 mirrors, with no direction. For the quiet design with a shared error,
differential evolution rescued up to 8 more mirrors in 100, so our
numbers for that case slightly underestimate what rescue can do. The re-designed
top layers are then ``deposited'' with the same error model as all
other layers --- the rescue does not get perfect layers. We test
$k = 0$ (no rescue), 2, 4, 6 and 8.

\subsection{Pass test and Monte Carlo}

A finished mirror \emph{passes} if its true transmissions land inside
the Advanced LIGO end-mirror windows \cite{E0900068}: between 4 and
6~ppm at 1064~nm, and between 1\% and 4\% at 532~nm. Thermal noise is
not part of the pass test; we report it separately, because the rescue
could in principle trade noise for yield.

For every combination of design, $\sigma$, measurement error and $k$,
we simulate 200 mirrors (400 in the fine scan) and report the fraction
that pass, with binomial uncertainties. The same random mirrors are
reused across all $k$ and all measurement levels, so differences
between settings come from the settings, not from luck. All code is
Python, with fixed random seeds, and is available at
\url{https://github.com/Tonmayyyyyyyyyyyy/LIGO-mirror-coating}.

% References you will need (fill in your .bib file):
% Venugopalan2024 - Optics Express 32(7) 11751 (2024), arXiv 2110.13437
% Hong2013       - Phys. Rev. D 87, 082001 (2013)
% E0900068       - LIGO document E0900068-v5
% MaterialsSource - arXiv 0912.0107 (title and authors to confirm)
% aLIGO2015     - Class. Quantum Grav. 32, 074001 (2015), arXiv 1411.4547
% Harry2002      - Class. Quantum Grav. 19, 897 (2002), first coating
%                  Brownian noise formula (cite in the introduction)
